\begin{problem}
    En la figura adjunta se representa una barra de \SI{20}{kg} con tres discos distintos en cada lado. Los discos de cada lado tienen una masa de \SI{5}{kg}, \SI{10}{kg} y \SI{15}{kg}.

    \begin{center}
    (imagen)
    \end{center}

    Si se quiere levantar la misma cantidad de kilogramos usando solo discos de \SI{10}{kg}, ¿cuántos discos se deben colocar a cada lado de la barra?

    \begin{enumerate}[label=\Alph*.]
        \item \num{3}
        \item \num{4}
        \item \num{5}
        \item \num{6}
    \end{enumerate}
\end{problem}
